\section{Extended MCMC Diagnostics}
\label{app:mcmc_diagnostics}

This appendix provides the extended sampler diagnostics requested
by reviewers for the non-convex Welsch generalised posterior. All
diagnostics are computed from the whale DGP at $N = 1000$ with
\texttt{severity=``severe''} across five contamination densities
($\{0\%, 1\%, 5\%, 10\%, 20\%\}$) and five seeds per density, for
a total of 25 runs. The default NUTS configuration (2 chains,
400 warmup, 800 samples) is used throughout.

\subsection{Convergence summary}
\label{app:mcmc:convergence}

\begin{table}[!t]
\centering
\small
\begin{tabular}{rcccc}
\toprule
density & $\hat{R}$ (max) & ESS (min) & IAC (max) & divergences \\
\midrule
0\%   & 1.01 & 518 & 3.1 & 0 \\
1\%   & 1.01 & 497 & 3.2 & 0 \\
5\%   & 1.01 & 496 & 3.2 & 0 \\
10\%  & 1.01 & 519 & 3.1 & 0 \\
20\%  & 1.01 & 531 & 3.0 & 0 \\
\bottomrule
\end{tabular}
\caption{Summary convergence diagnostics across contamination
densities (worst case over 5 seeds per density). $\hat{R}$:
split-$\hat{R}$ (Gelman--Rubin); ESS: minimum effective sample
size across parameters; IAC: maximum integrated
autocorrelation time in lags. All runs pass standard thresholds
($\hat{R} < 1.05$, ESS $> 200$) with substantial
margin.}
\label{tab:mcmc_diagnostics}
\end{table}

\subsection{Initialisation sensitivity}
\label{app:mcmc:init}
To test for multimodality or initialisation dependence, we refit
the whale DGP at 20\% density (the most challenging setting) with
three initialisation strategies: (i) default (prior-draw init),
(ii) OLS $\bar\beta_{\mathrm{OLS}}$ init, and (iii) random
overdispersed init ($\beta_0 \sim \mathcal{N}(0, 100)$). Across 3
seeds:

\begin{itemize}
\item Posterior means agree to within $0.02$ ($< 2\%$ of CI width)
  across all three strategies.
\item Posterior 95\% CI widths agree to within 5\%.
\item No bimodality observed in any bivariate scatter plot of
  $\beta$ draws.
\item Chain-specific means (2 chains $\times$ 3 strategies)
  cluster tightly around the pooled mean.
\end{itemize}

We conclude that in the $p \le 100$ regime tested, the
non-convexity of the Welsch $\rho_W$ does not induce practical
multimodality. The prior regularisation ($\beta \sim t_3(0,
\sigma_\beta^2)$) and the concentration of bulk residuals in the
locally convex region ($|r| < c/\sqrt{2} \approx 0.95$ at
$c = 1.34$) ensure that NUTS explores a single dominant mode.

\subsection{Autocorrelation analysis}
\label{app:mcmc:autocorrelation}
The integrated autocorrelation time (IAC) measures how many
sequential draws constitute one effective independent sample.
For all runs, IAC $\le 3.5$ lags. The IAC does not meaningfully
increase with contamination density (staying around $\sim$3.2
across all densities), indicating that the Welsch pseudo-likelihood's
non-convexity in the tail region does not impede NUTS exploration
at the densities tested.

\subsection{\texorpdfstring{Split-$\hat{R}$ on four chains}{Split-Rhat on four chains}}
\label{app:mcmc:split_rhat}
The main experiments use 2 chains. As a post-hoc robustness check,
we rerun 5 seeds of the whale DGP at 20\% density with 4 chains
(400 warmup, 800 samples each). Split-$\hat{R}$ (computed on the
4 half-chains, giving 8 segments) is $< 1.01$ for every parameter
at every seed, confirming that the 2-chain default is sufficient
for convergence assessment.
